\section{Recursive formulation: auditing the question, framework, and decomposition}
\label{sec:recursive-formulation}

Method evolution is incomplete if the method may change while the problem representation is
treated as fixed. For an unfamiliar research problem, the stated question, active representation,
decomposition, interface contracts, measurement model, and evaluator may themselves be the source
of error. We therefore treat these objects as versioned pursuit state. They may guide search, but
they receive no scientific authority merely because an executor selected them.

For an existing recursive fiber
\[
f=(o_f,q_f,\gamma_f,r_f,\operatorname{parent}(f)),
\]
define an audit projection
\[
\pi_{\mathrm{audit}}(f)=
(\mathbb Q_f,\Phi_f,\mathbb D_f,\mathbb I_f,\mathbb M_f,\mathbb V_f,
 \mathbb R_f,\chi_f),
\]
containing candidate question formulations, framework candidates, candidate decompositions,
interface contracts, measurement contracts, evaluator identity, residual-responsibility hypotheses
and recursive closure state. This is a projection of the existing research state rather than a
second authority architecture.

A recursive audit may recommend solving at the current representation, reframing the question,
challenging the framework, splitting or merging a provisional atom, repairing an interface,
revising a measurement contract, auditing the evaluator, running a discriminator, descending to a
child, ascending to an ancestor, stopping at a bounded cutoff, or returning
\texttt{CANNOT\_CHECK}. None of these actions changes scientific authority without the ordinary
certificate-producing gate.

Atomicity is scoped. A fiber is provisionally atomic only relative to a registered target, split
family, evaluator, evidence cutoff and decision consumer. A fresh discriminator that separates two
child regimes reopens that atomicity certificate without rewriting evidence accumulated under the
former representation.

Descendant failure may challenge an ancestor abstraction. Repeated raw failure is insufficient:
if several abstraction levels remain plausible, a decision-bearing discriminator is required.
When distinct local repair families fail and evidence implicates a parent assumption, the correct
action is upward reframing rather than indefinite downward patching. Prior descendants remain
immutable negative/process history under the superseded parent identity.

This calculus does not imply that one globally correct formulation exists or can always be found.
Closure remains target-, evaluator-, cutoff- and resource-relative.

\subsection{Formal obligations}
\label{sec:recursive-formulation-obligations}

The calculus carries six obligations, each falsifiable by a counterexample artifact rather than
by argument:

\begin{description}
\item[P-RF1 (Pursuit non-sovereignty)] Formulation change alone cannot increase scientific
  authority. Every audit action is a pursuit-state transition; authority still requires a
  certificate-producing evidence gate.
\item[P-RF2 (History preservation)] Superseding a question, framework or decomposition cannot
  delete old evidence or negative history; superseded identities remain addressable.
\item[P-RF3 (No post-hoc responsibility selection)] If multiple abstraction levels remain
  plausible, the system may not mutate all of them and call the winning level ``diagnosed'';
  a decision-bearing discriminator must run first.
\item[P-RF4 (Indexed atomicity)] Atomicity certificates are target-, split-family-, evaluator-
  and cutoff-indexed; no certificate subsumes another index, and there is no
  \texttt{ATOM\_PROVEN\_FOREVER} terminal.
\item[P-RF5 (Ancestor reopening)] Superseding an ancestor abstraction stales dependent descendant
  closure certificates while preserving their evidence identities.
\item[P-RF6 (Evaluator epoch closure)] An evaluator change closes the evaluation epoch; no
  cross-epoch comparison is offered. Method-scale audits reuse the strict Self-RAKL governance
  (Section~\ref{sec:upgrade}) and cannot bypass it.
\end{description}

\subsection{Executable known-world evaluation}
\label{sec:recursive-formulation-evaluation}

The obligations are executed against a frozen known-world benchmark
(\texttt{research/recursive\_framework\_audit\_v1/}): eleven defect families including a
well-posed control on which any reframe is harm, twelve hostile priority cases probing the
ordering of the audit's decision chain, and eight structural invariant checks instantiating
P-RF1--P-RF6. The production decision operator is differential-conformant to a vendored
handoff reference on every case. Conformance establishes that the control semantics are
executable and fail closed; it is instrument evidence, not utility evidence. Real open-ended
formulation utility is owned by the research-engine paper (Paper~VI), whose capstone separates
solution gain, responsible-coordinate localization, false-reframe harm, audit cost and
method-evolution gain.
